\section{The adequacy boundary and normalized public signatures}

Section~2 defined the raw sealed-extension calculus and the accounting invariant \(\mu\).
Those definitions are intentionally presentation-insensitive: they count only the normalized
opaque public signature exported by a sealed layer.  This section explains how an extension
declaration reaches that normal form.  It is the place where surface syntax, semantic
admissibility, opacity, and trace accounting meet.

The main point is a boundary rather than a slogan.  The paper proves its depth-two theorem
for the fixed raw calculus \(C_{\mathrm{ext}}\).  It does not prove a parser or elaborator
theorem for arbitrary Cubical Agda programs, and it does not assume that every cubical
foundation automatically presents its sealed extensions in exactly this way.  Instead, all
semantic readings pass through an \emph{adequacy bridge}: a package of preservation claims
connecting a surface or semantic extension language to the raw normalized signatures counted
in Section~2.

The section has four jobs.  First, it states the adequacy package that a surface language or
external cubical foundation must supply in order to inherit the results of the paper.  Second,
it gives a small concrete instance, \(\mathsf{SMod}\), showing that the boundary is not empty.
Third, it defines the presentation moves that normalization is allowed to forget.  Finally, it
proves that the resulting normalized public signatures preserve exactly the trace information
used later by the horn-computation, depth, and recurrence theorems.

A useful way to read this section is:
\[
\text{surface declaration}
\quad\leadsto\quad
\text{raw admissible declaration}
\quad\leadsto\quad
\text{sealed public interface}
\quad\leadsto\quad
\text{normalized public signature}.
\]
The depth theorem begins only after the last arrow.  This is what prevents the later
cardinality statements from depending on arbitrary choices of record packaging, telescope
bracketing, or explicit versus implicit derived fields.

\subsection{Semantic obligations and the adequacy package}

A sealed extension is not counted merely by looking at the symbols a programmer writes.  The
same public interface may be presented using a record, an iterated dependent sum, a bundled
operation, an unbundled telescope, or an explicit derived compatibility field.  Conversely,
small changes in surface notation may hide a real new primitive public obligation.  The
adequacy bridge identifies which fields are semantic payload, which fields are primitive
public trace, and which fields are derived before \(\mu\) is computed.

\begin{definition}[Semantic structural obligations]
Let \(E_{n+1}\) be a semantic sealed layer over an already sealed state \(B_n\).  A
\emph{semantic structural obligation} for \(E_{n+1}\) is an interface field required to
interpret the new payload of \(E_{n+1}\) against the opaque public interfaces exported by
previous layers.

Unary obligations specify the action of the new payload on one old public site.  Binary
obligations specify compatibility with an old comparison, path, transport, or equivalence
between public sites.  Higher horn-shaped obligations specify the remaining faces of a
higher structural compatibility problem generated by sealing.

User-authored higher payload is not, by itself, structural trace.  If a layer contributes a
higher inductive constructor, a path constructor, or another genuinely higher mathematical
operation, that data belongs to the payload part \(K_{n+1}\).  It becomes structural trace
only when it is a public compatibility field forced by sealing the new layer against old
opaque interfaces.
\end{definition}

The last sentence is one of the most important discipline rules of the paper.  The theorem
is not claiming that higher-dimensional mathematical content disappears.  It is claiming that
higher \emph{structural integration obligations} generated by sealing are represented as
cubical horn-extension problems and therefore do not survive as primitive public trace.

\begin{remark}[Adequacy package and status]
For a general cubical foundation \(\mathcal F\), or for a richer surface language intended to
compile into \(C_{\mathrm{ext}}\), the semantic reading of the theorems is parametric in an
adequacy package \(\mathcal A\).  The package contains the following preservation clauses.

\begin{enumerate}
\item \textbf{Raw soundness.}  Every well-typed surface payload, action, comparison,
transport, or horn-boundary declaration elaborates to the corresponding raw clause class of
\(C_{\mathrm{ext}}\).

\item \textbf{Raw completeness for the fixed sealing discipline.}  Every admissible semantic
structural obligation generated by the chosen sealing discipline is represented by one of the
raw clause classes of Section~2: payload, unary action, binary comparison, or structural horn
boundary.

\item \textbf{Opacity preservation.}  The surface notion of sealing agrees with the raw opaque
export.  After sealing, later layers may use the normalized public interface, but not hidden
transparent derivations except through explicitly exported fields or through the transparent
operations admitted by counting normalization.

\item \textbf{Support preservation.}  Normalization preserves the historical support of every
primitive trace field.  Repackaging may change notation, but it may not turn a field depending
on layers \(i,j\) into one depending on a different public history.

\item \textbf{Primitive/derived preservation.}  Presentation equivalence and normalization
preserve whether a field is primitive or transparently derived.  A field may be marked derived
only when the fixed derivability relation supplies a uniform construction from previously
exported lower data.

\item \textbf{Cardinality preservation.}  Payload size, primitive trace size, exported
interface cardinality, and minimal opaque cost \(\mu\) are preserved by the raw-to-normalized
signature bridge.

\item \textbf{Horn-computation preservation.}  Higher structural obligations that the surface
language treats as derived must elaborate to the structural horn-extension packages of
Sections~4 and~5.  In particular, their derived status must be witnessed by the cubical
operations used in the computational replacement theorem, not merely asserted by a tag.

\item \textbf{Coupling preservation.}  When a declaration advertises global action on an
active interface footprint, the elaborated raw declaration must cover exactly that footprint.
For the fully coupled recurrence theorem, this includes the per-site coverage assumptions of
Section~2.
\end{enumerate}

In this paper these clauses are proved or checked for the fixed raw extension calculus and
for the small \(\mathsf{SMod}\) surface instance below.  Transfer from arbitrary Cubical Agda
source code, or from a different cubical foundation, remains conditional on supplying the
corresponding adequacy package.  A proposed counterexample to the broad guiding conjecture
therefore has a precise target: it must either violate one of these bridge clauses or exhibit a
sealing-generated structural obligation that cannot be represented as a derived horn-extension
package in the target calculus.
\end{remark}

This formulation deliberately makes the broad transfer conjecture smaller and sharper.  The
paper proves a theorem inside \(C_{\mathrm{ext}}\); the adequacy package records exactly what
another system would have to preserve in order to inherit that theorem.

\subsection{A concrete adequacy instance: a global modality over a Tarski universe}

The adequacy boundary should not be an empty promise.  The accompanying development
therefore includes a small surface calculus, \(\mathsf{SMod}\), for adding a global modality
over an opaque old Tarski universe.  The example is intentionally modest.  Its role is not to
model every modality or every universe extension, but to show concretely how payload,
primitive trace, and derived horn trace are separated.

The old sealed interface contains a Tarski universe with function codes:
\[
\begin{array}{rcl}
U &:& \mathsf{Type},\\
\mathsf{El} &:& U \to \mathsf{Type},\\
\mathsf{Arr} &:& U \to U \to U,\\
\mathsf{ElArr} &:& (A\ B:U) \to
  \mathsf{El}(\mathsf{Arr}(A,B)) = (\mathsf{El}(A) \to \mathsf{El}(B)).
\end{array}
\]
A new modal declaration contributes payload
\[
\begin{array}{rcl}
\mathsf{Flat} &:& \mathsf{Type}\to\mathsf{Type},\\
\eta &:& (A:\mathsf{Type})\to A\to\mathsf{Flat}(A),
\end{array}
\]
and structural trace describing how \(\mathsf{Flat}\) acts on the old universe interface:
\[
\begin{array}{rcl}
\mathsf{FlatU} &:& U\to U,\\
\mathsf{flatEl} &:& (A:U)\to
  \mathsf{El}(\mathsf{FlatU}(A)) = \mathsf{Flat}(\mathsf{El}(A)),\\
\mathsf{flatCodeArr} &:& (A\ B:U)\to
  \mathsf{FlatU}(\mathsf{Arr}(A,B)) =
  \mathsf{Arr}(\mathsf{FlatU}(A),\mathsf{FlatU}(B)),\\
\mathsf{flatArr} &:& (A\ B:U)\to
  \mathsf{Flat}(\mathsf{El}(A)\to\mathsf{El}(B)) =
  (\mathsf{Flat}(\mathsf{El}(A))\to\mathsf{Flat}(\mathsf{El}(B))).
\end{array}
\]
The first two trace fields describe action and decoding.  The next two describe comparison
with the function-code structure: \(\mathsf{Flat}\) must interact coherently with both the code
\(\mathsf{Arr}(A,B)\) and the decoded function type.

For each \(A,B:U\), the old and new data generate two routes from
\[
  \mathsf{El}(\mathsf{FlatU}(\mathsf{Arr}(A,B)))
\]
to
\[
  \mathsf{Flat}(\mathsf{El}(A))\to\mathsf{Flat}(\mathsf{El}(B)).
\]
The left route first decodes \(\mathsf{FlatU}(\mathsf{Arr}(A,B))\) using
\(\mathsf{flatEl}\), transports the old decoding equation \(\mathsf{ElArr}\) through
\(\mathsf{Flat}\), and then applies \(\mathsf{flatArr}\).  The right route first uses
\(\mathsf{flatCodeArr}\) to compare function codes, decodes with \(\mathsf{ElArr}\) at the
lifted codes, and transports domain and codomain along \(\mathsf{flatEl}(A)\) and
\(\mathsf{flatEl}(B)\).  The comparison between these two routes is the generated horn field
\(\mathsf{flatArrCoh}\).

A surface program may present \(\mathsf{flatArrCoh}\) explicitly.  It may also omit it.  In
either case, the field is not primitive public trace: it is a structural horn whose lower faces
are the displayed unary and binary trace data, and Section~5 supplies the cubical filler.

\begin{center}
\begin{tabular}{llll}
\hline
Surface field & Role & Raw clause class & Counted by \(\mu\)? \\
\hline
\(\mathsf{Flat},\eta\) & payload & payload & no \\
\(\mathsf{FlatU},\mathsf{flatEl}\) & universe action/decoding & unary trace & yes \\
\(\mathsf{flatCodeArr},\mathsf{flatArr}\) & function comparison & binary trace & yes \\
\(\mathsf{flatArrCoh}\) & route coherence & derived horn & no \\
\hline
\end{tabular}
\end{center}

\begin{theorem}[Modal surface adequacy]
Let \(E\) be a well-typed declaration in \(\mathsf{SMod}\) over the opaque old universe
interface above.  Then \(\mathrm{Elab}(E)\) is an admissible raw
\(C_{\mathrm{ext}}\)-declaration.  Elaboration preserves payload fields, primitive trace
fields, historical support, and normalized public cost \(\mu\).  The generated
function-code compatibility horn elaborates to a derived structural horn clause.  Adding
that horn explicitly as a surface field does not change the normalized public signature or
\(\mu\).
\end{theorem}

\begin{proof}
The proof is syntax-directed.  Payload fields of \(\mathsf{SMod}\) elaborate to raw payload
clauses.  The fields \(\mathsf{FlatU}\) and \(\mathsf{flatEl}\) elaborate to unary action and
decoding clauses over the old universe interface.  The fields \(\mathsf{flatCodeArr}\) and
\(\mathsf{flatArr}\) elaborate to binary comparison clauses, because they compare the modal
action with already exported function-code and function-type structure.

The route comparison \(\mathsf{flatArrCoh}\) has the structural open-box shape analyzed in
Section~5.  Its boundary is assembled from the four trace fields above together with old
exported equalities and transparent transports.  If the surface declaration omits the horn,
elaboration inserts the canonical structural filler.  If the surface declaration supplies it
explicitly, contractibility of the theorem-facing horn-extension object identifies it with the
canonical filler at the level relevant to public normalization.  Hence the explicit and implicit
presentations normalize to canonically isomorphic public signatures with the same primitive
fields and the same value of \(\mu\).
\end{proof}

The mechanized instance is located in the modal surface modules:
\[
\begin{gathered}
\texttt{Surface/Modal/Syntax.agda},\quad
\texttt{Surface/Modal/Typing.agda},\quad
\texttt{Surface/Modal/Routes.agda},\\
\texttt{Surface/Modal/Elaboration.agda},\quad
\texttt{Surface/Modal/Normalization.agda},\quad
\texttt{Surface/Modal/Adequacy.agda}.
\end{gathered}
\]
The smoke test \texttt{Test/Surface/ModalAdequacySmoke.agda} checks the minimal
declaration, the explicit-horn presentation, and rebundling invariance of \(\mu\).

\subsection{Presentation steps and canonical normal forms}

The normalized public signature should not depend on how a programmer packages a public
interface.  A one-constructor record, a flattened telescope, and an equivalent iterated
\(\Sigma\)-type should present the same sealed public data.  Similarly, a derived alias should
not become primitive merely because it was written explicitly.  Presentation equivalence is
therefore generated by a small list of local moves.

\begin{definition}[Presentation-step generators]
Presentation equivalence for admissible declarations is generated by the following local
moves on normalized public telescopes:
\[
\begin{array}{ll}
\mathsf{reassocSigma} & \text{reassociate or flatten an iterated dependent sum;}\\
\mathsf{splitRecord} & \text{split a single-constructor record or shell into fields;}\\
\mathsf{bundleRecord} & \text{rebundle such a field telescope into a record or shell;}\\
\mathsf{curryPi} & \text{turn a product argument into an iterated function telescope;}\\
\mathsf{uncurryPi} & \text{perform the inverse uncurrying step;}\\
\mathsf{transportField} & \text{move a field along an already exported equality;}\\
\mathsf{transparentAliasInsert} & \text{insert a transparently definable alias field;}\\
\mathsf{transparentAliasDelete} & \text{delete such an alias field;}\\
\mathsf{duplicateDerivedDelete} & \text{delete a duplicate derived trace field.}
\end{array}
\]
Their reflexive, symmetric, and transitive closure is the relation
\[
  e \approx e'.
\]
The list is also the scope of the raw bridge.  It covers the raw constructors of
Section~2: payload telescopes, action clauses, comparison and transport clauses, algebraic
payload packages, export selections, and packaged horn-boundary clauses.  It does not include
arbitrary source-language rewrites or parser-level equivalence for Cubical Agda programs.
\end{definition}

The point of this definition is conservative.  Presentation equivalence forgets bureaucracy,
not semantic obligations.  A field may be removed only when it is an alias, a duplicate
derived field, or a different presentation of the same public typing judgment.

\begin{proposition}[Counting normalization is stable under presentation steps]
Each generator of presentation equivalence preserves the canonical counting normal form up to
telescope isomorphism, preserves primitive-versus-derived tags, and preserves historical
support.  Consequently, if \(e\approx e'\), then
\[
  N(e) \cong N(e'),
  \qquad
  \operatorname{Prim}(N_{\mathrm{tr}}(e)) \cong
  \operatorname{Prim}(N_{\mathrm{tr}}(e')),
  \qquad
  \mu(e)=\mu(e').
\]
\end{proposition}

\begin{proof}
Reassociation, record splitting, rebundling, currying, and uncurrying change only the shape of
the public telescope.  After flattening the \(\Sigma\)-spine and normalizing the
\(\Pi\)-telescope, both sides have the same atomic fields up to canonical renaming.

The transport move uses an equality already present in an earlier sealed public signature.
It therefore changes the displayed type of a field but not the irreducible prior public data
on which the field depends.  Transparent alias insertion and deletion add or remove only
fields whose inhabitants are uniformly synthesized from existing exported data.  Such fields
are tagged derived and do not contribute to primitive trace cost.  Duplicate-derived deletion
removes a second presentation of a field already classified as derived.

Thus every generator preserves primitive classes and their historical support.  Closure under
reflexivity, symmetry, and transitivity gives the result for \(e\approx e'\).
\end{proof}

In the accompanying development this is the theorem-facing contract of
\texttt{Metatheory/PresentationEquivalence.agda},
\texttt{Metatheory/TraceCostNormalForm.agda}, and
\texttt{Metatheory/MuInvariance.agda}.  The paper does not require a global confluence
theorem for arbitrary source syntax.  It requires uniqueness of the counted normal form for
this generated presentation relation, which is exactly the invariant used by the later
cardinality results.

\subsection{Elaboration, sealing, and public signatures}

We now connect admissible declarations to the candidate-and-obligation calculus of
Section~2.  The following theorem is the formal version of the pipeline displayed at the
start of the section.

\begin{theorem}[Elaboration soundness for sealed exports]
Every admissible surface extension declaration \(e\) over a state \(B_n\) determines:
\begin{enumerate}
\item an elaborated candidate
\[
  X(e) := (X_{\mathrm{core}}(e),G_{\mathrm{obl}}(e));
\]
\item a sealed layer \(L(e)\) exported through the opaque public interface selected by the
export component \(\sigma\);
\item a normalized public signature
\[
  N(e) := N(\mathrm{Elab}(e))
\]
obtained from that exported interface by counting normalization.
\end{enumerate}
In particular, the candidate calculus of Section~2 is the image of admissible surface
extensions under elaboration and sealing.
\end{theorem}

\begin{proof}
The payload telescope \(\Gamma_{\mathrm{new}}\) elaborates to new core declarations.  The
action telescope \(\Gamma_{\mathrm{act}}\) and comparison telescope \(\Gamma_{\mathrm{cmp}}\)
elaborate to the fields witnessing interaction with the active interface of \(B_n\).  The
admissibility judgment says exactly that these fields type-check and discharge the induced
obligations of the declared coupling mode.  The export selector \(\sigma\) then determines the
opaque public telescope of the sealed layer.  Applying counting normalization to that public
telescope yields \(N(e)\).  These data are precisely the candidate and sealed layer used by
Section~2.
\end{proof}

\begin{theorem}[Adequacy of normalized public signatures]
Let \(e\) be an admissible surface extension over \(B_n\), and let \(L(e)\) be its sealed
export.  Then:
\begin{enumerate}
\item every field of \(N(e)\) comes from a field of the sealed public interface of \(L(e)\);
counting normalization introduces no new primitive public datum;
\item every field of \(N(e)\) is either a payload schema contributed by
\(\Gamma_{\mathrm{new}}\) or a trace schema arising from a resolved structural obligation of
\(X(e)\);
\item every irreducible payload or trace schema visible after sealing occurs as a field of
\(N(e)\);
\item the quantities
\[
  P(X(e)),\quad S(L(e)),\quad \Delta(X(e)\mid B_n),\quad \mu(e),\quad
  \operatorname{ObSig}_k(X(e)),\quad \operatorname{Ix}_k(X(e)),\quad
  \operatorname{Real}_k(X(e))
\]
are exactly the corresponding payload family, public signature, abstract latency, minimal
opaque cost, obligation signature, primitive index set, and realized obligation object computed
from \(N(e)\) and its trace part.
\end{enumerate}
\end{theorem}

\begin{proof}
The first claim follows from the definition of sealing and counting normalization: the
normalization procedure flattens the exported telescope, quotients presentation artifacts, and
discards uniformly synthesized fields, but it does not invent primitive public data.

The second claim is the surface-language form of the payload/trace decomposition of
Section~2.  After elaboration and sealing, a visible field is either newly introduced
mathematical content of the layer or public evidence that the new content interacts correctly
with earlier opaque interfaces.

For the third claim, observe that counting normalization removes only presentation artifacts
and transparently reconstructible fields.  If an irreducible field is visible after sealing, it is
not uniformly reconstructible from the remaining exported data by the admitted transparent
operations.  It therefore survives into \(N(e)\).

The final claim is bookkeeping.  Section~2 defines payload size, public schema set,
integration latency, minimal opaque cost, obligation signatures, primitive index sets, and
realized obligation objects from normalized public signatures and their induced trace
obligations.  Since \(N(e)\) is the canonical normalized signature of the sealed export, these
abstract quantities agree with the corresponding surface quantities after elaboration.
\end{proof}

The theorem explains why the later sections can stop mentioning surface syntax.  Once a
declaration has crossed the adequacy bridge, all remaining claims are claims about normalized
public signatures and their induced structural obligations.

\begin{proposition}[Necessity of irreducible public fields]
Let \(e\) be an admissible surface extension over \(B_n\), and let \(\phi\) be an irreducible
field of \(N(e)\).  If \(e'\approx e\), then \(N(e')\) contains a canonically corresponding
irreducible field \(\phi'\).  Equivalently, no presentation-equivalent declaration can delete an
irreducible payload or primitive trace schema from the sealed public signature.
\end{proposition}

\begin{proof}
Presentation equivalence identifies elaborated sealed extensions only up to definitional
equality, canonical telescope isomorphism, and transparent internal normalization.  If the
irreducible field \(\phi\) were absent from \(N(e')\), then either the public typing judgment of
\(e'\) would genuinely omit a field present in the public typing judgment of \(e\), contradicting
\(e'\approx e\), or \(\phi\) would be uniformly reconstructible from the remaining exported data
by the transparent operations admitted by presentation equivalence.  The latter contradicts
irreducibility.  Hence a canonical corresponding field must remain.
\end{proof}

\begin{proposition}[Canonical normalized presentations]
If \(e\approx e'\), then \(N(e)\) and \(N(e')\) are canonically isomorphic, and their trace
parts are canonically isomorphic.  Consequently every presentation-equivalence class admits a
canonical normalized public signature up to isomorphism, and
\[
  \mu(e)=\left|\operatorname{Prim}(N_{\mathrm{tr}}(e))\right|.
\]
\end{proposition}

\begin{proof}
Presentation equivalence identifies elaborated sealed extensions whenever they induce the
same public typing judgment up to the generated presentation steps.  Counting normal form
forgets exactly those presentation differences and nothing more.  Therefore \(N(e)\) and
\(N(e')\) are canonically isomorphic.  The payload/trace decomposition is part of the sealed
public judgment, so the trace subfamilies correspond under this isomorphism.  By the
necessity proposition, presentation equivalence cannot remove an irreducible primitive trace
field.  Hence the minimum defining \(\mu\) is attained by the canonical normalized trace
signature itself, and its value is the cardinality of the primitive part of that trace signature.
\end{proof}

\subsection{Computational replacement on canonical trace presentations}

The previous results say that packaging changes do not alter \(\mu\).  The next theorem says
something stronger and more important: an explicitly written higher structural field can
vanish from primitive trace cost when a cubical term reconstructs it from its boundary.
This is the formal expression of the phrase ``derivedness is a theorem, not a tag.''

The theorem below does not itself construct the cubical term.  That construction is supplied
by the horn-to-open-box and open-box computation results of Sections~4 and~5.  What the
theorem proves is that, once such a term exists in the fixed derivability relation, the field is
not allowed to remain in the \(\mu\)-minimal primitive public signature.

\begin{theorem}[Computational replacement on canonical trace presentations]
Let \(X\) be a candidate over \(B_n\), let \(k\) be a history depth, and let \(S\) be the
canonical normalized trace signature attached to the depth-\(k\) trace family of \(X\).  Let
\(P\) and \(Q\) be two presented trace telescopes whose counting normal forms are \(S\), and
let \(\phi\) be a presented trace field of \(P\).

Suppose there is a lower-arity boundary subfamily \(\partial\phi\) and a cubically synthesized
term
\[
  t_\phi(\partial\phi)
\]
built only from \(\mathsf{hcomp}\), \(\mathsf{transp}\), \(\mathsf{hfill}\), definitional
equality, and the transparent structural operations already quotiented by counting
normalization, such that after passage to \(S\) the field corresponding to \(\phi\) is
judgmentally reconstructed from lower-arity boundary data.  Then:
\begin{enumerate}
\item \(P\) and \(Q\) are presentation-equivalent representations of the same canonical trace
signature;
\item in the \(\mu\)-minimal part of the canonical normalized trace signature represented by
\(S\), the field corresponding to \(\phi\) is absent as an irreducible primitive trace field;
\item every other irreducible primitive trace field of \(P\) corresponds canonically to one of
\(Q\).
\end{enumerate}
\end{theorem}

\begin{proof}
Both \(P\) and \(Q\) normalize to \(S\), so they are two presentations of one canonical public
trace interface.  The hypothesis gives more than a label for \(\phi\): it gives a uniform term
\(t_\phi(\partial\phi)\) in the fixed transparent derivability relation.  That term reconstructs
the normalized image of \(\phi\) from lower-arity boundary data using only cubical Kan
operations and presentation-transparent structural operations.

Consequently the normalized image of \(\phi\) is tagged transparently derived rather than
primitive.  The \(\mu\)-minimal primitive signature retains only irreducible primitive trace
fields, so \(\phi\) is excluded from that primitive part.  Since primitive/derived tags and
support are preserved under passage between presentations of \(S\), every other irreducible
primitive field of \(P\) corresponds canonically to one of \(Q\).  Thus computational
replacement changes the presentation of the trace telescope, not the primitive public content
measured by \(\mu\).
\end{proof}

This theorem is the bridge between the cubical argument and the accounting argument.  The
open-box theorem will show that the relevant higher structural fields have terms of the form
\(t_\phi(\partial\phi)\).  The computational replacement theorem then turns that cubical fact
into the public-signature fact that those fields do not contribute primitive trace cost.

\begin{remark}[Exact scope of the mechanized fixed raw bridge]
The formal bridge in the accompanying Agda development is a three-part package for the
fixed raw extension calculus.  It is not an additional metatheory of all cubical foundations
and not a parser for arbitrary Cubical Agda source.

\begin{enumerate}
\item Raw admissible declarations are connected, via elaboration and sealing, to canonical
normalized signatures by
\texttt{Metatheory/RawStructuralSyntax.agda},
\texttt{Metatheory/RawStructuralTyping.agda}, and
\texttt{Metatheory/SurfaceNormalizationBridge.agda}.

\item Presentation-equivalent canonical trace normal forms preserve support, primitive cost,
and \(\mu\) by the explicit presentation-step constructors in
\texttt{Metatheory/PresentationEquivalence.agda} and
\texttt{Metatheory/MuInvariance.agda}.

\item The resulting canonical normalized signatures are matched with the counted-interface
quantities by the normalized-signature/counting bridge, and well-typed raw structural clauses
are identified with the horn-generated trace image by
\texttt{Metatheory/SurfaceToHornImage.agda}.
\end{enumerate}

In particular, the computational replacement theorem is formalized at the level of canonical
trace presentations.  The bridge modules connect admissible raw declarations to that
interface and to the counted abstract interface.  What remains deliberately outside the
present theorem is a general elaboration theorem for the whole Cubical Agda language.
\end{remark}

\subsection{The trace-cost dictionary}

The recurrence theorem later uses the symbols \(\mu_k\), \(\kappa_k\), \(\Delta\), and
\(S_{\mathrm{bas}}(L_k)\).  Those symbols live at slightly different levels: surface
declarations, normalized trace signatures, abstract obligation families, and primitive basis
classes.  The next two results record that, after the adequacy bridge, they are the same
cardinalities read through different lenses.

\begin{corollary}[Minimal opaque cost agrees with abstract coherence cost]
For every admissible surface extension declaration \(e\) over \(B_n\),
\[
  \mu(e)
  = \left|\operatorname{Prim}(N_{\mathrm{tr}}(e))\right|
  = \Delta(X(e)\mid B_n).
\]
If \(e_k\) is the declaration whose sealing exports \(L_k\), if
\(\mu_k:=\mu(e_k)\), and if \(S_{\mathrm{bas}}(L_k)\) is any basis family of \(L_k\), then
\[
  |S_{\mathrm{bas}}(L_k)| = \kappa_k + \mu_k.
\]
\end{corollary}

\begin{proof}
By adequacy of normalized public signatures, the trace part of \(N(e)\) records exactly the
irreducible public trace schemas arising from resolved obligations of \(X(e)\).  By canonical
normalization, \(\mu(e)\) is the cardinality of the primitive part of that trace signature.  By
the active-basis accounting of Section~2, these primitive trace schemas are exactly the basis
sites counted by \(\Delta(X(e)\mid B_n)\).  The final formula is the payload/trace decomposition
of the sealed public basis: the primitive basis of \(L_k\) consists of payload classes of size
\(\kappa_k\) together with primitive trace classes of size \(\mu_k\).
\end{proof}

\begin{lemma}[Trace-cost dictionary]
Let \(e_k\) be an admissible fully coupled declaration whose sealing exports \(L_k\).  Then
the following quantities are the same cardinality, read at their respective levels:
\[
  \mu(e_k)
  = \left|\operatorname{Prim}(N_{\mathrm{tr}}(e_k))\right|
  = \Delta(X(e_k)\mid B_{k-1})
  = |\operatorname{Str}(L_k)|.
\]
Moreover, the primitive basis of the exported layer decomposes as
\[
  |S_{\mathrm{bas}}(L_k)|
  = |S_{\mathrm{core}}(L_k)| + |\operatorname{Str}(L_k)|
  = \kappa_k + \mu_k.
\]
Thus the recurrence later uses a cardinality identity obtained from normalized public
signatures, not an additional semantic growth assumption.
\end{lemma}

\begin{proof}
The first equality is the definition of minimal opaque cost on the canonical normalized trace
presentation.  The second is the previous corollary.  The third reads the same primitive
trace classes as the structural part of the sealed export.  The second displayed formula is
the primitive-class restriction of the same payload/trace decomposition, with
\(|S_{\mathrm{core}}(L_k)|=\kappa_k\).  Invariance under presentation replacement follows from
presentation-step stability.
\end{proof}

\subsection{Factorization-complete public trace export}

The depth theorem is an arity theorem: it says that primitive structural trace stabilizes at
binary comparison.  The recurrence theorem needs one more ingredient.  A binary comparison
could still point far back in history, so arity depth two does not automatically imply a
two-layer chronological window.  The following definition isolates the extra export property
needed to reuse recent public trace without reopening hidden derivations.

\begin{definition}[Factorization-complete public trace export]
Let \(L_k\) be a sealed layer with \(k\geq 2\).  Its public trace export is
\emph{factorization-complete} when, after choosing any admissible declaration whose sealing
exports \(L_k\) and reading its canonical normalized trace signature, the following data are
available for the next remote-comparison step:
\begin{enumerate}
\item the adjacent unary bridge from \(L_k\) to \(L_{k-1}\);
\item the binary comparison trace witnessing the composite comparison from \(L_k\) to
\(L_{k-2}\) through \(L_{k-1}\);
\item the degenerate or structural faces needed to assemble the corresponding open
\(3\)-box, generated transparently by definitional equality, transport along already exported
equalities, and the structural operations already quotiented by counting normalization.
\end{enumerate}
Equivalently, the next horn-extension problem involving \(L_k\) can be posed using only
exported public trace fields and transparent structural operations, without reopening the
hidden derivation that discharged the original sealing obligations.
\end{definition}

\begin{remark}[Why the transparency boundary is not arbitrary]
The phrase ``generated transparently'' is part of the theorem-facing boundary.  It is not a
license to declare inconvenient fields derived.  A field is transparent only when the fixed
derivability relation exhibits a uniform construction from already exported lower data, using
definitional equality, canonical telescope isomorphisms, transport along exported equalities,
or the cubical Kan operations admitted by the raw bridge.

Changing that derivability relation would change the measured primitive signature.  This is
why the paper treats the transparency boundary as part of adequacy rather than as informal
notation.  The two-layer window depends on a checkable claim: a remote comparison field
must be computed by the horn term supplied later by the structural open-box theorem and the
computational replacement theorem.  If a proposed foundation cannot supply such a uniform
term, the remote field remains primitive and the foundation does not instantiate the
chronological factorization theorem.
\end{remark}

\begin{theorem}[Sealed layers export factorization-complete trace data]
Every sealed layer \(L_k\) with \(k\geq 2\) in the fixed fully coupled cubical extension calculus
has factorization-complete public trace export.
\end{theorem}

\begin{proof}
Choose an admissible declaration \(e_k\) whose sealing exports \(L_k\).  By the
payload/trace decomposition, the public export of \(L_k\) consists of new payload together
with trace.  By adequacy of normalized public signatures, every irreducible trace schema
visible after sealing appears in the canonical normalized trace signature
\(N_{\mathrm{tr}}(e_k)\).

Because \(L_k\) is sealed as a globally acting layer against the active interface, its resolved
trace obligations include the adjacent action on \(L_{k-1}\) and the binary comparison against
the already exported comparison between \(L_{k-1}\) and \(L_{k-2}\).  Hence the corresponding
unary bridge and binary comparison trace occur as public fields of the normalized trace
signature.

The remaining faces needed for the next open \(3\)-box are not extra primitive obligations.
They are degenerate composites, reindexings, and transports along already exported equalities,
so they are generated by the transparent structural operations admitted in counting
normalization.  Presentation equivalence cannot delete the irreducible unary or binary trace
fields just identified; it can only rewrite them into canonically isomorphic normal form.
Therefore the factorization data are exported stably at the level of canonical normalized
public signatures and can be reused without reopening the hidden sealing derivation.
\end{proof}

\begin{remark}[Use of the bridge in later sections]
Sections~4 and~5 invoke surface syntax only through the adequacy bridge just described.  Once
an admissible declaration has been identified with its canonical normalized signature and the
corresponding counted-interface data, the later horn, depth, factorization, and recurrence
theorems are applications of the abstract candidate-and-obligation calculus.
\end{remark}
